\chapter{Weighted curvature and smoothing of periodic graphs}\label{sec:graphical-weighted-curvature}
Retain the conventions of Section \ref{sec:graphical-flow-angle}. The cancellation behind the estimates of \cite{EH89}, Lemma 4.1, Corollary 4.2 and Proposition 4.4, printed pp.~459--461, has an exact tensor-square form. It avoids differentiation of $|A|$ at a zero of curvature.

\begin{proposition}[An exact square identity for slope-weighted curvature]\label{prop:graphical-curvature-square}
On a smooth graphical mean curvature flow in flat ambient space put
\[
 P=v^2|A|^2,\qquad
 \mathcal D=\mathcal P+2v^{-1}\nabla v\cdot\nabla.
\]
Then
\[
 \mathcal DP=-2|\nabla(vA)|^2.
\]
The norm uses the evolving induced metric, and the identity remains smooth at $A=0$.
\end{proposition}
\begin{proof}
Let $a=|A|^2$. Proposition \ref{prop:mcf-curvature-evolution} gives $\mathcal Pa=-2|\nabla A|^2+2a^2$, and Proposition \ref{prop:graphical-angle-evolution} gives $\mathcal P(v^2)=-2av^2-6|\nabla v|^2$. Expanding the product yields
\[
 \mathcal P(av^2)=-2v^2|\nabla A|^2-6a|\nabla v|^2
                 -4v\langle\nabla a,\nabla v\rangle.
\]
The quartic curvature terms cancel exactly. The added drift contributes $2v\langle\nabla a,\nabla v\rangle+4a|\nabla v|^2$. Thus
\[
 \mathcal DP=-2v^2|\nabla A|^2-2a|\nabla v|^2
              -2v\langle\nabla a,\nabla v\rangle.
\]
Since $\nabla a=2\langle A,\nabla A\rangle$, this is the expansion of $-2|\nabla(vA)|^2$. The norm evolution cited above already includes the evolving-metric terms; no fixed-metric differentiation of a covariant curvature tensor has been substituted. Its proof and the product calculation also apply in dimension one. This refines the nonpositive inequality used in \cite{EH89}, Lemma 4.1.
\end{proof}

\begin{corollary}[Weighted curvature cannot grow on a periodic flow]\label{cor:periodic-weighted-curvature-bound}
On a smooth periodic graph over $[s,T)$,
\[
 \max_{M_t}v^2|A|^2\leq\max_{M_s}v^2|A|^2.
\]
This is an estimate in terms of the initial curvature as well as the initial slope.
\end{corollary}
\begin{proof}
Proposition \ref{prop:graphical-curvature-square} gives $\mathcal DP\leq0$. Lemma \ref{lem:graphical-material-drift} converts it into a scalar inequality with a smooth drift on each compact periodic slab. The same first-maximum argument as in Proposition \ref{prop:periodic-graph-slope-bound} bounds its supremum by the initial one. In particular no bound on the derivatives of $|A|$ is required to apply this inequality at a point with zero curvature. Compare \cite{EH89}, Corollary 4.2.
\end{proof}

\begin{lemma}[A quantity that eliminates the initial curvature bound]\label{lem:graphical-smoothing-quantity}
For a fixed starting time $s$, define
\[
 Z=v^2+2(t-s)v^2|A|^2.
\]
At every smooth graphical point,
\[
 \mathcal DZ=-2|\nabla v|^2-4(t-s)|\nabla(vA)|^2.
\]
In particular this quantity is a subsolution for $t\geq s$.
\end{lemma}
\begin{proof}
The power equation and the drift give
$\mathcal D(v^2)=-2|A|^2v^2-2|\nabla v|^2=-2P-2|\nabla v|^2$.
The time derivative of $2(t-s)P$ supplies $2P$, and Proposition \ref{prop:graphical-curvature-square} supplies the remaining negative square. Addition proves the formula. The coefficient two in $Z$ is exactly what cancels the undifferentiated $P$ term. This is the calculation underlying the case $m=0$ of \cite{EH89}, Proposition 4.4.
\end{proof}

\begin{theorem}[Curvature smoothing from slope alone]\label{thm:periodic-graph-curvature-smoothing}
For a smooth periodic graphical flow on $[s,T)$, let $V_s=\max_{M_s}\sqrt{1+|Du|^2}$. Then for $s<t<T$,
\[
 \bigl(1+2(t-s)|A(x,t)|^2\bigr)v(x,t)^2\leq V_s^2,
 \qquad
 |A(x,t)|^2\leq\frac{V_s^2-1}{2(t-s)}.
\]
The first estimate retains more information than the second. It depends on the starting slope but not its curvature. This is the compact periodic form of the base smoothing estimate in \cite{EH89}, Proposition 4.4, p.~461.
\end{theorem}
\begin{proof}
Lemma \ref{lem:graphical-smoothing-quantity} and the compact drift maximum principle imply $Z(x,t)\leq\max Z(\cdot,s)=V_s^2$. This is the first inequality. Since $v\geq1$, it implies $2(t-s)|A|^2\leq V_s^2/v^2-1\leq V_s^2-1$, proving the second. The proof is on any smooth compact slab, so the conclusion holds at every time before $T$. It is an a priori estimate for a flow already known to exist; a continuation or global-existence theorem requires its own parabolic input.
\end{proof}

\begin{corollary}[Uniformly bounded ancient periodic slopes are rigid]\label{cor:ancient-periodic-graph-rigidity}
Suppose $u:\mathbb T^n\times(-\infty,T)\to\mathbb R$ is a smooth periodic graphical flow and
$\sup_{x,t}\sqrt{1+|Du(x,t)|^2}<\infty$.
Then $u$ is a single constant, independent of space and time. No bound on its height over the entire ancient interval is required.
\end{corollary}
\begin{proof}
Fix $t<T$ and apply Theorem \ref{thm:periodic-graph-curvature-smoothing} with any $s<t$. The numerator is uniformly bounded, while $t-s\to\infty$ as $s\to-\infty$. Hence $A(\cdot,t)=0$. On a graph, $h_{ij}=-u_{ij}/v$, so the Euclidean Hessian of the periodic lift vanishes. Its gradient is a constant vector; periodicity along every vector of a full-rank lattice forces that vector to be zero. Thus $u$ is spatially constant at each time. The graphical equation then gives $u_t=0$. The uniform slope bound is used for all past times, not merely separately on each time slice.
\end{proof}

\begin{example}[Why the grim reaper does not violate ancient rigidity]\label{ex:grim-slope-curvature-identity}
For the grim-reaper graph $u(x,t)=-\log\cos x+t$ on $(-\pi/2,\pi/2)$, and for its products with Euclidean factors,
\[
 v=\sec x,\qquad |A|=\cos x,\qquad v^2|A|^2=1.
\]
The tensor $vA$ is parallel along each time slice. Thus it realizes the zero right-hand side of Proposition \ref{prop:graphical-curvature-square}, even though its curvature is not zero. Its unbounded slope and its strip domain exclude both the periodic smoothing conclusion with finite $V_s$ and Corollary \ref{cor:ancient-periodic-graph-rigidity}.
\end{example}
\begin{proof}
The derivatives are $u_x=\tan x$ and $u_{xx}=\sec^2x$. With upward normal and $A=D\nu$, the signed principal curvature of the curve is $-u_{xx}/(1+u_x^2)^{3/2}=-\cos x$. Multiplication by $v$ makes it $-1$. In dimension one this is a constant multiple of the identity tensor, hence parallel. For its product with a flat factor, the product connection preserves the curve and flat tangent factors, and $vA$ is minus the projection onto the curve factor, again parallel. The displayed identities follow. The slope tends to infinity at both strip edges, so the finite starting supremum used in the smoothing proof is unavailable. This example checks the exact square without extending a compact maximum principle to an unbounded domain.
\end{proof}
